\subsection{Transitivity of Reduced Norms and Traces}

PROPOSITION 7. --- Let L be a maximal commutative semisimple subalgebra of A and a an element of L. We have

\[
\mathrm{Pcrd} _ {\mathrm{A/K}} (a; \mathrm{X}) = \mathrm{Pc} _ {\mathrm{L/K}} (a; \mathrm{X}),\tag{52}
\]

\[
\mathrm{Trd} _ {\mathrm{A/K}} (a) = \mathrm{Tr} _ {\mathrm{L/K}} (a),\tag{53}
\]

\[
\mathrm{Nrd} _ {\mathrm{A/K}} (a) = \mathrm{N} _ {\mathrm{L/K}} (a).\tag{54}
\]

By Proposition 3 of VIII, p.~262, the L-modules A and \(L^{n}\) are isomorphic; we therefore have the relation

\[
\mathrm{Pc} _ {\mathrm{A/K}} (a; \mathrm{X}) = \mathrm{Pc} _ {\mathrm{L/K}} (a; \mathrm{X}) ^ {n}.
\]

Since the polynomial \(\mathrm{Pc}_{\mathrm{L}/\mathrm{K}}(a;\mathrm{X})\) is monic, it is therefore equal to the reduced characteristic polynomial \(\mathrm{Pcrd}_{\mathrm{A}/\mathrm{K}}(a;\mathrm{X})\) (VIII, p.~339, Lemma 2); this gives formula (52). By comparing the coefficients of \(X^{n-1}\) (resp. the constant terms) on each side of (52), we obtain formula (53) (resp. (54)).

COROLLARY. --- Let D be a field of finite degree over K with center K. Let a be an element of K and L a maximal commutative subfield of D containing a. We have

\[
\begin{array}{c} \mathrm {Pcrd_ {D / K}} (a; \mathrm{X}) = \mathrm {Pc_ {L / K}} (a; \mathrm{X}), \\ \mathrm{Tr} _ {\mathrm{D/K}} (a) = \mathrm{Tr} _ {\mathrm{L/K}} (a), \\ \mathrm{Nrd} _ {\mathrm{D/K}} (a) = \mathrm{N} _ {\mathrm{L/K}} (a). \end{array}\tag{55}
\]

Indeed, a maximal commutative subfield L of D is a maximal commutative semisimple subalgebra of D by Corollary 2 of VIII, p.~265.

PROPOSITION 8. --- Let B be a simple subalgebra of A. Denote the center of B by L and the commutant of B in A by B'. Then B' is a central simple algebra over the field L; we denote its reduced degree by r. For every element b of B, we have the relations

\[
\mathrm{Pcrd} _ {\mathrm{A/K}} (b; \mathrm{X}) = \mathrm{N} _ {\mathrm{L[X]/K[X]}} (\mathrm{Pcrd} _ {\mathrm{B/L}} (b; \mathrm{X})) ^ {r},\tag{56}
\]

\[
\mathrm{Trd} _ {\mathrm{A/K}} (b) = r \mathrm{Tr} _ {\mathrm{L/K}} (\mathrm{Trd} _ {\mathrm{B/K}} (b)),\tag{57}
\]

\[
\mathrm{Nrd} _ {\mathrm{A/K}} (b) = \mathrm{N} _ {\mathrm{L/K}} (\mathrm{Nrd} _ {\mathrm{B/L}} (b)) ^ {r}.\tag{58}
\]

Lemma 6. --- Let \(\mathbf{K}'\) be a commutative algebra of finite degree \(d\) over \(\mathbf{K}\) and

\[
\mathrm{P} (\mathrm{X}) = \mathrm{X} ^ {s} + a _ {1} \mathrm{X} ^ {s - 1} + \dots + a _ {s}
\]

a monic polynomial with coefficients in \(K'\) . Then the polynomial \(Q = N_{K'[X]/K[X]}(P)\) in \(K[X]\) is monic of degree sd, the coefficient of \(X^{sd-1}\) in \(Q(X)\) is equal to \(\mathrm{Tr}_{\mathrm{K}'/\mathrm{K}}(a_1)\) , and the constant term of Q is \(N_{K'/K}(a_s)\) .

We denote the \(K^{\prime}\) -algebra \(\mathrm{K}^{\prime}[\mathrm{T}]/( \mathrm{P}(\mathrm{T}))\) by \(K^{\prime\prime}\) and the canonical class of T in \(K^{\prime\prime}\) by t. The sequence \((1, t, \ldots, t^{s-1})\) is a basis of \(K^{\prime\prime}\) over \(K^{\prime}\) , and the matrix of multiplication by t in this basis is of the form

\[
\tau = \left( \begin{array}{c c c c c} 0 & 0 & \dots & 0 & - a _ {s} \\ 1 & 0 & \dots & 0 & - a _ {s - 1} \\ 0 & 1 & \dots & 0 & - a _ {s - 2} \\ . & . & \dots & . & . \\ . & . & \dots & . & . \\ 0 & . & \dots & 1 & - a _ {1} \end{array} \right).\tag{59}
\]

The determinant of \(XI_{n}-\tau\) is calculated by induction on s, by expanding along the first row. We obtain \(\det(\mathrm{X}I_{n}-\tau)=\mathrm{P}(\mathrm{X})\) . In other words, we have \(\mathrm{P}(\mathrm{X})=\mathrm{P}\mathrm{c}_{\mathrm{K}^{\prime\prime}/\mathrm{K}^{\prime}}(t;\mathrm{X})\) . In particular, \(\mathrm{Tr}_{\mathrm{K}^{\prime\prime}/\mathrm{K}^{\prime}}(t)=-a_{1}\) and \(\mathrm{N}_{\mathrm{K}^{\prime\prime}/\mathrm{K}^{\prime}}(t)=(-1)^{s}a_{s}\) . By the transitivity formula (III, §9, No.~4, p.~548, Corollary), we have

\[
\begin{array}{r} \mathrm{Tr} _ {\mathrm{K} ^ {\prime \prime} / \mathrm{K}} (t) = - \mathrm{Tr} _ {\mathrm{K} ^ {\prime} / \mathrm{K}} (a _ {1}), \quad \mathrm{N} _ {\mathrm{K} ^ {\prime \prime} / \mathrm{K}} (t) = (- 1) ^ {s d} \mathrm{N} _ {\mathrm{K} ^ {\prime} / \mathrm{K}} (a _ {s}), \\ \mathrm{Q} (\mathrm{X}) = \mathrm{Pc} _ {\mathrm{K} ^ {\prime \prime} / \mathrm{K}} (t; \mathrm{X}). \end{array}
\]

On the other hand, \([K'' : K] = [K'' : K'][K' : K] = sd\) , so \(Q(X)\) is a monic polynomial of degree sd. Lemma 6 follows.

Let us prove Proposition 8. Since the ring B is simple, its center L is a field (VIII, p.~121, Corollary 1). By Theorem 5 of VIII, p.~259, the commutant \(B'\) of B in A is a simple ring with center L, and we have the equality \([A:K]=[B:K][B':K]\) . We denote the reduced degree of \(B'\) over L by r, that of B over L by s, and the degree of L over K by d.~We have

\[
[ \mathrm{A}: \mathrm{K} ] = n ^ {2}, [ \mathrm{B} ^ {\prime}: \mathrm{K} ] = r ^ {2} d, [ \mathrm{B}: \mathrm{K} ] = s ^ {2} d,
\]

and therefore \(n^{2}=r^{2}s^{2}d^{2}\) , that is, n=rsd.

Let \(b\) be an element of B, and let \(\mathrm{P}(\mathrm{X})\) be its reduced characteristic polynomial over the L-algebra B; it is monic of degree \(s\) . By Lemma 6, the polynomial \(\mathrm{Q} = \mathrm{N}_{\mathrm{L}[\mathrm{X}] / \mathrm{K}[\mathrm{X}]}(\mathrm{P})\) is monic of degree \(sd\) . The polynomial \(\mathrm{R} = \mathrm{Q}^r\) is therefore monic of degree \(rsd = n\) . Again by Lemma 6, the coefficient of \(\mathrm{X}^{n - 1}\) in \(\mathrm{R}(\mathrm{X})\) is equal to \(-r\operatorname{Tr}_{\mathrm{L} / \mathrm{K}}(\operatorname{Trd}_{\mathrm{B} / \mathrm{L}}(b))\) , and the constant term of \(\mathrm{R}(\mathrm{X})\) is \((\mathrm{N}_{\mathrm{L} / \mathrm{K}}((-1)^s\mathrm{Nrd}_{\mathrm{B} / \mathrm{L}}(b)))^r = (-1)^n\mathrm{N}_{\mathrm{L} / \mathrm{K}}(\mathrm{Nrd}_{\mathrm{B} / \mathrm{L}}(b))^r\) .

Since \([A:K]=r^{2}d[B:K]\) , the left B-module A is free of rank \(r^{2}d\) (VIII, p.~124, Proposition 5). We therefore have

\[
\mathrm{Pc} _ {\mathrm{A/K}} (b; \mathrm{X}) = \mathrm{Pc} _ {\mathrm{B/K}} (b; \mathrm{X}) ^ {d r ^ {2}}.\tag{60}
\]

By the corollary of III, §9, No.~4, p.~548, we have

\[
\mathrm{Pc} _ {\mathrm{B/K}} (b; \mathrm{X}) = \mathrm{N} _ {\mathrm{L[X] / K[X]}} (\mathrm{Pc} _ {\mathrm{B/L}} (b; \mathrm{X})),\tag{61}
\]

and since \(P(X)\) is the reduced characteristic polynomial of b over the L-algebra B, we have

\[
\mathrm{Pc} _ {\mathrm{B/L}} (b; \mathrm{X}) = \mathrm{P(X)} ^ {s}.\tag{62}
\]

Finally, by formulas (60)---(62) and the definition of R(X), we have

\[
\mathrm{Pc} _ {\mathrm{A/K}} (b; \mathrm{X}) = \mathrm{N} _ {\mathrm{L[X]/K[X]}} (\mathrm{P(X)}) ^ {d r ^ {2} s} = \mathrm{Q(X)} ^ {d r ^ {2} s} = \mathrm{R(X)} ^ {r s d} = \mathrm{R(X)} ^ {n}\tag{63}
\]

so \(\mathrm{R}(\mathrm{X})\) is the reduced characteristic polynomial of b over the K-algebra A.

We have proved formula (56). Formulas (57) and (58) follow immediately from formula (56) and Lemma 6 because the coefficient of \(\mathrm{X}^{n - 1}\) in \(\mathrm{Pcrd}_{\mathrm{A / K}}(b;\mathrm{X})\) is equal to \(-\mathrm{Trd}_{\mathrm{A / K}}(b)\) and the constant term is \((-1)^n\mathrm{Nrd}_{\mathrm{A / K}}(b)\) .

